\subtituloA{Arquitectura de referencia del navegador}

\nota{Para el examen}{
La arquitectura de referencia es el conjunto de componentes básicos comunes a la mayoría de los navegadores modernos.
}

\subtituloB{Componentes principales (Parte 1)}

\rista{
\lista{
  **Interfaz de usuario:** Actúa como interfaz entre el usuario y el motor, gestionando barras de herramientas, progreso, preferencias y comunicación con el sistema operativo. |
  **Motor del buscador:** Carga la dirección (URL/URI), soporta navegación básica (anterior/siguiente/recargar) y gestiona alertas de JavaScript. |
  **Motor de renderizado:** Presenta visualmente el recurso, interpretando HTML, XML, CSS y contenido embebido (audio/vídeo/imágenes). |
  **Subsistema de comunicaciones:** Implementa protocolos de transferencia de archivos y documentos de internet (ej. HTTP/HTTPS).
}
}

\subtituloB{Componentes principales (Parte 2)}

\rista{
\lista{
  **Intérprete de JavaScript:** Analiza y ejecuta código JavaScript, respondiendo a eventos (ratón, teclado) o código intercalado. |
  **Parser XML:** Carga en memoria una representación en árbol (DOM) del documento HTML para facilitar y agilizar el acceso a sus elementos. |
  **Componente de visualización:** Ofrece funcionalidades específicas para la presentación del contenido HTML. |
  **Subsistema de persistencia de datos:** Almacena historiales, sesiones de usuario, preferencias, certificados de seguridad y cookies.
}
}

\nota{Para el examen}{
¿Cuál es la función del parser XML en el contexto de la visualización web?

\textbf{Respuesta:} Cargar una representación en árbol (DOM) del documento HTML en memoria para agilizar el acceso a sus elementos.
}
